\subsection{The finite Hopf-algebra count used by the free-\(O\) producer}
\label{subsec:faithful-hopf-count}

This subsection translates the finite algebra argument in
\texttt{PrimitiveDerivationPairing}, together with its five elementary
producers.  Throughout, \(k\) is an algebraically closed field of
characteristic a prime \(p\), and \(A\) is a finite-dimensional,
commutative and cocommutative Hopf \(k\)-algebra.  Write \(\Delta\),
\(\epsilon\), and \(\iota\) for its comultiplication, counit, and
antipode.  Its finite dual \(A^\vee=\operatorname{Hom}_k(A,k)\) is
endowed with convolution.  In particular
\[
 (uv)(a)=(u\otimes v)(\Delta a),\qquad
 1_{A^\vee}=\epsilon .
\]
An \(\epsilon\)-derivation is a \(k\)-linear map \(w:A\to k\) with
\(w(ab)=w(a)\epsilon(b)+\epsilon(a)w(b)\).  An element \(f\in A\) is
primitive when \(\Delta f=f\otimes1+1\otimes f\).

\begin{lemma}[Two finite-algebra choices]\label{lem:faithful-hopf-choices}
Let \(R\) be a finite-dimensional commutative \(k\)-algebra.
\begin{enumerate}
\item If a \(k\)-subspace \(U\subset R\) has the property that every
\(k\)-algebra homomorphism \(x:R\to k\) is nonzero on some element of
\(U\), then \(U\) contains a unit.
\item If \(V\ne0\) is a \(k\)-subspace of \(R\), stable under
\(f\mapsto f^p\), then it contains a nonzero \(f\) for which
\(f^p=0\) or \(f^p=f\).
\end{enumerate}
\end{lemma}
\begin{proof}
The ring \(R\) is Artinian.  Every maximal-ideal residue field is a
finite extension of \(k\), hence is \(k\).  In (1), the intersection
of \(U\) with each maximal ideal is therefore a proper subspace.
There are finitely many maximal ideals, and \(k\) is infinite.
A finite union of proper subspaces cannot cover \(U\).  An element
outside that union belongs to no maximal ideal, and is a unit.

For (2), use strong induction on \(\dim_k(fR)\), beginning with any
nonzero \(f\in V\).  If \(f\) is nilpotent, take the last nonzero term
in \(f,f^p,f^{p^2},\ldots\).  Otherwise \(f\) is outside some prime
ideal, which is maximal.  The corresponding \(x:R\to k\) has
\(x(f)\ne0\).  Set \(f'=x(f)^{-1}f\) and \(g=(f')^p-f'\).
Then \(g\in V\) and \(x(g)=0\).  If \(g=0\), use \(f'\).
If \(g\ne0\), then \(g=fq\) for some \(q\in R\), so \(gR\subset fR\).
This inclusion is strict: equality would express \(f=ga\), contrary
to \(x(f)\ne0=x(g)\).  Apply the induction hypothesis to \(g\).
\end{proof}

\begin{lemma}[Powers of a nonzero primitive]\label{lem:faithful-primitive-powers}
If \(0\ne f\in A\) is primitive, the \(p\) elements
\(1,f,\ldots,f^{p-1}\) are linearly independent.
\end{lemma}
\begin{proof}
Applying the counit to one tensor factor of the primitive identity
gives \(\epsilon(f)=0\).  Thus \(f\) and \(1\) are linearly independent.
Choose a linear functional \(y:A\to k\) with \(y(1)=0\) and \(y(f)=1\).
Let \(y^{*m}\) denote its \(m\)-fold convolution power, with
\(y^{*0}=\epsilon\).  The binomial formula gives
\[
 \Delta(f^i)=\sum_{l=0}^i {i\choose l}f^l\otimes f^{i-l}.
\]
Induction on \(m\) consequently gives
\[
 y^{*m}(f^i)=0\quad(i<m),\qquad
 y^{*m}(f^m)=m!.
\]
Indeed, in the convolution step all terms except \(l=1\) vanish:
the \(l=0\) term uses \(y(1)=0\), and for \(l\ge2\) the other
convolution factor has exponent less than \(m\).
Given a linear relation among the displayed \(p\) powers, apply
\(y^{*m}\) successively with \(m=p-1,p-2,\ldots,0\).
Since \(m!\ne0\) for \(m<p\), all coefficients vanish.
\end{proof}

We record the two derivation computations in the precise form used
below.  For a derivation \(D:R\to R\), put \(C=\ker D\).

\begin{lemma}[Nilpotent derivation coordinates]
\label{lem:faithful-nilpotent-derivation}
Suppose \(R\) is a finite-dimensional commutative \(k\)-algebra,
\(D^p=0\), and \(D^{p-1}t\) is a unit for some \(t\in R\).
Then there is \(z\in R\) with \(Dz=1\), and the map
\[
 C^p\longrightarrow R,\qquad (c_0,\ldots,c_{p-1})
       \longmapsto\sum_{i=0}^{p-1}c_i z^i
 \tag{HC1}\label{eq:faithful-nilpotent-coordinates}
\]
is a \(k\)-linear isomorphism.  In particular
\(\dim_kR=p\dim_k C\).
For every \(k\)-algebra map \(e:R\to k\), an \(e\)-derivation
\(u:R\to k\) which kills \(C\) and \(z\) is zero; moreover \(e(Dz)=1\).
\end{lemma}
\begin{proof}
Put \(w=D^{p-1}t\).  Since \(Dw=0\), also \(D(w^{-1})=0\), and
\[
 z=w^{-1}D^{p-2}t
\]
has \(Dz=1\).  This formula includes \(p=2\).
For \(c\in C\),
\[
 D^n(cz^i)=i(i-1)\cdots(i-n+1)c z^{i-n},
\]
where a term with \(n>i\) is zero.  The coefficient formula follows
by induction from Leibniz's rule.  It proves injectivity of
\eqref{eq:faithful-nilpotent-coordinates}, by applying \(D^{p-1}\)
to a relation and descending on the largest exponent.
For surjectivity, prove by induction on \(j\le p\) that
\(D^ja=0\) implies an expression with exponents \(<j\).
At the induction step, \(b=D^ja\) lies in \(C\); subtract
\((j!)^{-1}bz^j\), then apply the assertion for \(j\).
The factorials in use have \(j<p\).  Taking \(j=p\) proves the claim.
Finally \(u(z^i)=0\) by induction using its Leibniz rule, and
\(u(c_i)=0\); the expression in \eqref{eq:faithful-nilpotent-coordinates}
therefore gives \(u=0\).
\end{proof}

\begin{lemma}[Multiplicative derivation coordinates]
\label{lem:faithful-multiplicative-derivation}
Suppose \(R\) is a finite-dimensional commutative \(k\)-algebra,
\(D^p=D\), and \(z\) is a unit with \(Dz=z\).  Then
\[
 C^p\xrightarrow{\ \sim\ }R,\qquad
 (c_0,\ldots,c_{p-1})\longmapsto\sum_{i=0}^{p-1}c_i z^i .
 \tag{HC2}\label{eq:faithful-multiplicative-coordinates}
\]
Thus \(\dim_kR=p\dim_k C\).  For any \(k\)-algebra map \(e:R\to k\),
an \(e\)-derivation \(u\) killing \(C\) and \(z\) is zero, and
\(e(Dz)=e(z)\ne0\).
\end{lemma}
\begin{proof}
The polynomial \(X^p-X\) has simple roots, exactly the elements of
\(\mathbf F_p\subset k\).  It annihilates \(D\), so \(R\) is the
direct sum of its \(D\)-eigenspaces with eigenvalues
\(0,1,\ldots,p-1\).  Leibniz's rule gives
\(D(z^i)=i z^i\) and \(D(z^{-i})=-i z^{-i}\).
Multiplication by \(z^i\) therefore identifies \(C\) with the
eigenspace of eigenvalue \(i\): its inverse is multiplication by
\(z^{-i}\).  This proves \eqref{eq:faithful-multiplicative-coordinates}.
The last assertion follows from the same Leibniz calculation as
in Lemma~\ref{lem:faithful-nilpotent-derivation}.
\end{proof}

\begin{lemma}[A single primitive Hopf quotient]
\label{lem:faithful-single-primitive-quotient}
Let \(0\ne f\in A\) be primitive, with \(f^p=0\) or \(f^p=f\).
Then
\[
 \dim_k A=p\dim_k(A/(f)).
 \tag{HC3}\label{eq:faithful-single-primitive-dimension}
\]
Furthermore, the primitive elements in the ideal \((f)\) are
exactly the scalar multiples of \(f\).
\end{lemma}
\begin{proof}
On the commutative convolution algebra \(R=A^\vee\), define
\[
 (Du)(a)=u(fa).
\]
The primitive identity for \(f\) says exactly that \(D\) is a
derivation.  Its iterates satisfy
\((D^nu)(a)=u(f^na)\).  Let \(e:R\to k\) be evaluation at \(1_A\);
then \(e(Du)=u(f)\).

We first obtain one of the two coordinate lemmas.
If \(f^p=0\), then \(D^p=0\).  The element \(f^{p-1}\) is nonzero
by Lemma~\ref{lem:faithful-primitive-powers}.  Every algebra map
\(x:A^\vee\to k\) is evaluation at an invertible element \(g\in A\).
Here is the exact finite-dual argument: finite-dimensional
biduality represents \(x\) by \(g\); multiplicativity and unitality
give \(\Delta g=g\otimes g\) and \(\epsilon(g)=1\); the antipode
identity gives \(g\iota(g)=1\).
Choose \(u\in A^\vee\) with \(u(f^{p-1}g)\ne0\).  Then
\(x(D^{p-1}u)\ne0\).
Lemma~\ref{lem:faithful-hopf-choices}(1), applied to
\(\operatorname{im}D^{p-1}\), supplies a unit in that image.
Lemma~\ref{lem:faithful-nilpotent-derivation} applies.

If \(f^p=f\), then \(D^p=D\).  Set
\(b=\sum_{i=1}^{p-1}f^i\).
The equality \(fb=b\) follows from \(f^p=f\), and \(b\ne0\) follows
from Lemma~\ref{lem:faithful-primitive-powers}.
For an algebra map \(x\) represented by an invertible \(g\) as above,
choose \(u\) with \(u(bg)\ne0\).
The functional \(w(a)=u(ba)\) satisfies \(Dw=w\), and \(x(w)\ne0\).
Lemma~\ref{lem:faithful-hopf-choices}(1) now supplies a unit in
the eigenvalue-one subspace of \(D\).
Lemma~\ref{lem:faithful-multiplicative-derivation} applies.

In either case we have
\[
 \dim_kR=p\dim_k\ker D
 \tag{HC4}\label{eq:faithful-dual-rank}
\]
and an element \(z\in R\) such that \(e(Dz)\ne0\), with the following
property: an \(e\)-derivation killing \(\ker D\) and \(z\) is zero.
The kernel of \(D\) is the annihilator of the ideal \((f)\) under
the pairing \(A^\vee\times A\to k\).  Finite-dimensional duality
identifies its dimension with \(\dim_k A/(f)\), proving
\eqref{eq:faithful-single-primitive-dimension}.

Let \(g=bf\) be primitive.  Since \(z(f)=e(Dz)\ne0\), set
\(a=z(g)/z(f)\) and define
\[
 u(w)=w(g)-a\,w(f)\qquad(w\in R).
\]
The primitive identities for \(g\) and \(f\) make \(u\) an
\(e\)-derivation.  If \(Dw=0\), then \(w\) annihilates \((f)\),
so \(u(w)=0\); also \(u(z)=0\) by the choice of \(a\).
The property just proved forces \(u=0\).
The perfect pairing then gives \(g=af\).
Conversely, scalar multiples of \(f\) are primitive.
\end{proof}

\begin{proposition}[Primitive--tangent divisibility]
\label{prop:faithful-primitive-tangent-count}
Let \(P\subset A\) be a subspace of primitive elements, and
\(W\subset A^\vee\) a subspace of \(\epsilon\)-derivations, with
\[
 w(f)=0\qquad(w\in W,\ f\in P).
\]
Then
\[
 p^{\dim_kP+\dim_kW}\ \mid\ \dim_k A .
 \tag{HC5}\label{eq:faithful-primitive-tangent-count}
\]
No reducedness, connectedness, or perfect pairing between these
two subspaces is assumed.
\end{proposition}
\begin{proof}
Use strong induction on \(\dim_k A\).  We first carry out the step
when \(P\ne0\).  Enlarge it to
\[
 Q=\{f\in A:\ f\text{ is primitive and }w(f)=0
                    \text{ for every }w\in W\}.
\]
This is a nonzero subspace.  It is stable under \(p\)-th powers:
the primitive identity is preserved by Frobenius, and the
\(\epsilon\)-derivation rule gives \(w(f^p)=0\).
Lemma~\ref{lem:faithful-hopf-choices}(2) supplies
\(0\ne f\in Q\) with \(f^p=0\) or \(f^p=f\).

The ideal \(I=(f)\) is a Hopf ideal.  Indeed \(\epsilon(f)=0\);
the primitive identity becomes zero after applying the quotient
to both tensor factors; and \(\iota(f)=-f\), by the antipode
identity.  These assertions extend from \(f\) to its ideal by
multiplicativity.  Hence \(B=A/I\) is again a finite commutative
cocommutative Hopf algebra.  Write \(\pi:A\to B\).
By Lemma~\ref{lem:faithful-single-primitive-quotient},
\[
 \dim_kA=p\dim_kB,\qquad
 \ker(\pi|_Q)=kf .
 \tag{HC6}\label{eq:faithful-hopf-induction-dim}
\]
In particular \(\dim_kB<\dim_kA\) and
\(\dim_k\pi(Q)=\dim_kQ-1\).

Every \(w\in W\) kills \(I\): for \(af\in I\), its Leibniz rule
and \(w(f)=\epsilon(f)=0\) give \(w(af)=0\).
It therefore descends uniquely to an \(\epsilon_B\)-derivation on
\(B\).  Let \(W_B\subset B^\vee\) be the space of these descended
maps.  Pullback by the surjection \(\pi\) identifies \(W_B\) with
\(W\), so their dimensions agree.  The primitive subspace
\(P_B=\pi(Q)\) is annihilated by \(W_B\).
The induction hypothesis gives
\[
 p^{\dim_kQ-1+\dim_kW}\mid\dim_kB.
\]
Multiplying by \(p\) and using \eqref{eq:faithful-hopf-induction-dim}
gives the bound with \(Q\) in place of \(P\).
Since \(P\subset Q\), this implies \eqref{eq:faithful-primitive-tangent-count}.

If \(P=0=W\), the assertion is immediate.  If \(P=0\) but \(W\ne0\),
pass to the finite dual Hopf algebra \(A^\vee\).  The defining
Leibniz equality says precisely that the image of \(W\) in
\(A^\vee\) consists of primitive elements: this follows by testing
the dual comultiplication against arbitrary \(a\otimes b\).
Apply the nonzero-primitive induction step just proved to that
subspace and the zero tangent subspace.  Its induction hypotheses
are available because all smaller quotient dimensions are smaller
than \(\dim_kA^\vee=\dim_kA\).
This yields \(p^{\dim_kW}\mid\dim_kA^\vee=\dim_kA\), as required.
\end{proof}

\paragraph{Exact source ledger and remaining interface.}
All source links below refer to the frozen baseline, and each
listed file was compared byte for byte with its translation input.
The five elementary producers are
\href{https://github.com/ChristianKudeba/alpoge-rh-verification/blob/alpoge-source-7cbf7a0c/src/Development/AutomorphicSpectrum/Proofs/GeometryProofs/FrobeniusStableSubspaceElement.lean}
{\texttt{FrobeniusStableSubspaceElement}}
(blob \texttt{15fb69f4ae27a6ab7224fa1bf79bce80e0675cc2}),
\href{https://github.com/ChristianKudeba/alpoge-rh-verification/blob/alpoge-source-7cbf7a0c/src/Development/AutomorphicSpectrum/Proofs/GeometryProofs/SubspaceContainsUnit.lean}
{\texttt{SubspaceContainsUnit}}
(\texttt{8704c238b26ef7343b77b18851d75a59aee66050}),
\href{https://github.com/ChristianKudeba/alpoge-rh-verification/blob/alpoge-source-7cbf7a0c/src/Development/AutomorphicSpectrum/Proofs/GeometryProofs/NilpotentDerivationRank.lean}
{\texttt{NilpotentDerivationRank}}
(\texttt{ea20fda1a391f68e58f8727f9f03d43b523e6939}),
\href{https://github.com/ChristianKudeba/alpoge-rh-verification/blob/alpoge-source-7cbf7a0c/src/Development/AutomorphicSpectrum/Proofs/GeometryProofs/MultiplicativeDerivationRank.lean}
{\texttt{MultiplicativeDerivationRank}}
(\texttt{979dd7f8e24e2ac0afd291416a0ccd017b285dd1}),
and
\href{https://github.com/ChristianKudeba/alpoge-rh-verification/blob/alpoge-source-7cbf7a0c/src/Development/AutomorphicSpectrum/Proofs/GeometryProofs/PrimitiveElementPowersIndep.lean}
{\texttt{PrimitiveElementPowersIndep}}
(\texttt{51bb9dcde6060dee77bc4d8e4c835a0abc59281b}).
The assembly is
\href{https://github.com/ChristianKudeba/alpoge-rh-verification/blob/alpoge-source-7cbf7a0c/src/Development/AutomorphicSpectrum/Proofs/GeometryProofs/PrimitiveDerivationPairing.lean}
{\texttt{PrimitiveDerivationPairing}}
(\texttt{27d2edd4e342ad3e38cf97e7a3adc841abb3cd48}),
specifically \texttt{key}, \texttt{stepP},
\texttt{norm\_index\_count\_bound}, and
\texttt{primitive\_tangent\_dvd\_finrank}.
The proofs above retain its primitive quotient and dual-derivation
induction; they do not substitute a general structure theorem for
finite group schemes.

The next geometric interface is separate: for the finite kernel
of multiplication by \(a\), one must construct the injected
primitive and tangent subspaces, and identify their pairing with
the actual mixed coefficient of the adjoint cubical frame.
Proposition~\ref{prop:faithful-primitive-tangent-count} proves the
finite algebra count once those maps are supplied.  It does not
itself identify the germ's doubled tangent space with a
Dieudonn\'e module or establish the free-\(O\) conclusion.
